\begin{pdfdiagram}[x=1cm,y=1cm]
  \node[hw label, font=\small\bfseries] at (2.75,4.25) {SLC：2 态，1 条读参考线};
  \node[hw label, font=\small\bfseries] at (8.15,4.25) {TLC：8 态，7 条读参考线};
  % 两个独立面板避免把不同示意比例误读成同一连续电压轴。
  \fill[BookCodeBg] (0.6,0.75) rectangle (4.9,3.75);
  \draw[BookRule, line width=0.45pt] (0.6,0.75) rectangle (4.9,3.75);
  \fill[BookCodeBg] (5.45,0.75) rectangle (10.85,3.75);
  \draw[BookRule, line width=0.45pt] (5.45,0.75) rectangle (10.85,3.75);
  \draw[BookRule, line width=0.6pt, ->] (0.6,0.75) -- (5.10,0.75);
  \draw[BookRule, line width=0.6pt, ->] (5.45,0.75) -- (11.05,0.75)
    node[right, hw label] {$V_{th}$ (V)};
  \draw[BookRule, line width=0.6pt, ->] (0.6,0.75) -- (0.6,4.0);
  \node[hw label, rotate=90] at (0.18,2.25) {单元数};
  \draw[BookRule, line width=0.6pt, ->] (5.45,0.75) -- (5.45,4.0);
  \begin{scope}
    \clip (0.6,0.75) rectangle (4.9,3.75);
    \draw[BookTeal, line width=0.95pt]
      (0.95,0.92) .. controls (1.35,3.45) and (1.85,3.45) .. (2.25,0.92);
    \draw[BookTeal, line width=0.95pt]
      (3.25,0.92) .. controls (3.65,3.45) and (4.15,3.45) .. (4.55,0.92);
    \draw[BookRed, line width=0.6pt, dashed] (2.75,0.75) -- (2.75,3.55);
  \end{scope}
  \node[hw label, anchor=north] at (1.60,0.67) {E};
  \node[hw label, anchor=north] at (3.90,0.67) {P};
  \node[BookRed, font=\tiny\bfseries] at (2.75,3.55) {$V_{ref}$};
  \draw[BookAmber, line width=0.55pt, <->] (1.60,1.25) -- (3.90,1.25);
  \node[hw label, fill=white, inner sep=1pt] at (2.75,1.25) {状态间距约 5 V};

  \begin{scope}
    \clip (5.45,0.75) rectangle (10.85,3.75);
    \foreach \x in {5.72,6.36,7.00,7.64,8.28,8.92,9.56,10.20} {
      \draw[BookAccent, line width=0.72pt]
        (\x,0.92) .. controls ({\x+0.17},2.88) and ({\x+0.31},2.88) .. ({\x+0.48},0.92);
    }
    \foreach \x in {6.24,6.88,7.52,8.16,8.80,9.44,10.08} {
      \draw[BookRed, line width=0.42pt, dashed] (\x,0.75) -- (\x,3.30);
    }
  \end{scope}
  \foreach \x/\s in {5.96/E,6.60/P1,7.24/P2,7.88/P3,8.52/P4,9.16/P5,9.80/P6,10.44/P7} {
    \node[hw label, font=\tiny, anchor=north] at (\x,0.67) {\s};
  }
  \node[BookRed, font=\tiny\bfseries] at (8.16,3.48) {7 个 $V_{ref}$ 把相邻状态分开};
  \draw[BookAmber, line width=0.55pt, <->] (7.88,1.30) -- (8.52,1.30);
  \node[hw label, fill=white, inner sep=1pt, anchor=south] at (8.20,1.35) {约 0.8 V};
  \node[hw label, text width=9.8cm] at (5.72,-0.48)
    {读法：钟形曲线表示同一编程状态在大量单元上的 $V_{th}$ 分布；红色虚线是判决边界。状态越多，相邻分布越近，工艺波动、读干扰和保持退化越容易让分布尾部越过判决线。};
\end{pdfdiagram}
